\chapter{Natural numbers, counting, and induction}\label{ch:naturals}
\begin{readergoal}See how counting can be built from basic rules, rather than accepted as an unexplained operation.\end{readergoal}
\section{Counting precedes measurement}
When we say there were seven action potentials, we count discrete events. Counting does not require time units, voltage calibration, or a model of firing probability. The natural numbers are the first formal setting for such statements. We adopt $\mathbb N_0=\{0,1,2,\ldots\}$, reserving $\mathbb N_+=\{1,2,\ldots\}$. Authors differ about including zero, so every technical treatment should make its convention explicit.

\section{Successors and the Peano viewpoint}
Imagine a set $N$ with a distinguished element $0$ and successor operation $S:N\to N$. The basic idea is that zero is not the successor of anything; distinct inputs have distinct successors; and any subset containing zero and closed under the successor operation contains all of $N$. The final clause is the induction principle. These structural properties establish a unique natural-number system up to isomorphism, under the usual axiomatic background \cite{MendelsonNumbers}.

Define addition recursively by $n+0=n$ and $n+S(m)=S(n+m)$. Define multiplication through $n\cdot0=0$ and $n\cdot S(m)=n\cdot m+n$. These are \emph{definitions}, while commutativity, associativity, and distributivity are derived properties. We ordinarily write $n+1$ rather than $S(n)$ after connecting the symbols.

\section{Induction as a way to prove infinite families of claims}
Suppose we wish to prove $\sum_{k=1}^n k=n(n+1)/2$ for all positive integers $n$. The base case $n=1$ holds. Assume it is true for $n$. Then
\[\sum_{k=1}^{n+1}k=\frac{n(n+1)}2+(n+1)=\frac{(n+1)(n+2)}2.\]
Therefore the claim holds for $n+1$; induction establishes every case. Induction is a proof principle, not an empirical inference from observing a finite number of examples.

\section{Finite sets and combinatorics}
Two finite sets have the same cardinality if there exists a bijection between them. The number of length-$n$ binary vectors is $2^n$; the number of ways to choose $k$ distinct objects from $n$ without order is $\binom nk=n!/[k!(n-k)!]$. If 20 channels are inspected in all unordered pairs, the number of pairs is $\binom{20}{2}=190$. This count matters when reporting a family of connectivity tests.

\begin{workedexample}[Discrete time windows]
A recording sampled at $f_s=1000$ samples/s for 10 seconds has $N=f_sT=10000$ samples if we choose the index set $\{0,\ldots,9999\}$. The sample at index $n$ corresponds to $t_n=n/f_s$, so the last recorded time is $9.999$ s, not 10 s. Distinguish the \emph{duration of acquisition} from the \emph{timestamp of its final sample}.
\end{workedexample}
\begin{exercises}
\item Prove by induction that $1+3+\cdots+(2n-1)=n^2$.
\item Count unordered neuron pairs from $N=80$ recorded units.
\item Show that a finite sequence of $N$ entries is a function on $\{0,\ldots,N-1\}$.
\item Explain why mathematical induction and statistical generalization from trials are different operations.
\end{exercises}